\section{Integrating the period map over the Ran space}
\label{sec:ran}

The period map of Section~\ref{sec:unipotent} becomes an isomorphism after
allowing additional marked points and integrating over them.  This is the
remaining ingredient in the pairing identity \eqref{eq:T1}.  We first prove
an incidence calculation for principal parts, then apply it to unipotent
cochains and to the PBW filtration for restriction from $G$ to $N$.
Throughout, the level divisor is supported on the original tuple $T$.
Additional markings carry zero level.  When making underlying-complex
calculations, we apply the forgetful functor $b$ of
Section~\ref{sec:loop} in the original tuple variables.  The added
markings retain their relative D-module operations.

For an intermediate tuple $U\supseteq T$, insertion and the map
\eqref{eq:T10}, or its restriction version \eqref{eq:T15}, give
\begin{equation}\label{eq:T16}
 l_N[-\delta_N]\otimes
 \int_{V\supseteq U} B^0_V\bigl(\operatorname{ins}_{U,V}M\bigr)
 \longrightarrow \operatorname{Per}^{\mathrm{lv}}_U(M)[-2\delta].
\end{equation}
Here the target is identified using propagation \eqref{eq:T3} and the Ran
trace.  The symbol $\operatorname{Per}$ denotes the $N$-period for an
$N$-input and the $G$-period for a $G$-input.  After applying $b$ in the old variables, we
may equivalently use the source
$l^{\mathrm{lv}}[\delta]\otimes B^{\mathrm{lv}}$ and unshifted periods,
by \eqref{eq:T9}.

\begin{proposition}\label{prop:ran-integrated-period}
The map \eqref{eq:T16} is an isomorphism for unipotent Harish--Chandra
inputs at every intermediate tuple $U\supseteq T$.  With $U=T$, it is
also an isomorphism for every $G$ Harish--Chandra input whenever the
level divisor is a nonzero admissible thickening of the entire tuple $T$.
These isomorphisms are
continuous, linear over the D-module category of the original parameter
scheme, and compatible with insertion, deeper levels, and the parameter
comparisons of Sections~\ref{sec:localization} and~\ref{sec:unipotent}.
\end{proposition}

The base-change convention in this statement is essential.  Integration is
relative to the old base.  After applying $b$, crystal base change computes it
by the same Ran colimits and relative de Rham operations in the added
marking schemes.  The proofs below consequently retain the crystal
structures in these added variables.

\subsection{Principal parts and incidence}

Let $J_0$ be a vector bundle on the curve over an affine base.  For an
additional tuple $V$, let $P_V(J_0)=\operatorname{PP}(J_0,V)$ be the full
relative principal parts: the filtered union over all pole orders, modulo
regular sections, pushed from the divisor support.  We equip this object
with its principal-part crystal in the varying tuple, in the
angle-bracket dualizing normalization.  On the curve set
\[
 \mathcal I_X(J_0)=
 \bigl(\mathscr D_X\otimes_{\mathcal O_X}
               (J_0\otimes\Omega_X^{-1})\bigr)^{\langle X\rangle}.
\]
All these objects and operations are relative to the affine base.  Write
$i$ for the closed incidence $z\in V$ in the product of the curve with the
tuple scheme, and $p_V$ for projection from that product to the tuple
scheme.

\begin{lemma}\label{lem:ran-principal-parts}
There is a natural identification
\begin{equation}\label{eq:T17}
 P_V(J_0)\simeq
 p_{V,*}i_*i^!\bigl(\mathcal I_X(J_0)
                         \boxtimes\omega_{\mathrm{tuple}}\bigr).
\end{equation}
Under this identification, forgetting supports is the canonical map
$P_V(J_0)\to R\Gamma(X,J_0)[1]$.  For every finite list of vector bundles
$J_1,\ldots,J_r$ and every old tuple $U$, these maps induce an isomorphism
\begin{equation}\label{eq:T18}
 \int_{V\supseteq U}\mathop{\bigotimes}\limits_{j=1}^r{}^{!} P_V(J_j)
 \xrightarrow{\ \sim\ }
 \bigotimes_{j=1}^r R\Gamma(X,J_j)[1].
\end{equation}
The assertion includes $r=0$, with the tensor units understood, and is
compatible with permutations, shifts, and symmetric powers.
\end{lemma}

\begin{proof}
Apply the localization triangle to the complement of the incidence.  This
complement is the complement of the Cartier marking divisor, with
multiplicities allowed.  In our pushforward convention the relative-curve
de Rham complex of $\mathcal I_X(J_0)$ computes $J_0[1]$.  Indeed, passing
from a left induced module to a right module tensors by $\Omega_X$;
the factor $\Omega_X^{-1}$ in the definition cancels it, and the
angle-bracket normalization supplies the smooth-curve integration shift,
exactly as in \eqref{eq:T14}.  Localizing differential operators gives the
same calculation on the open complement.  The complex with supports is
therefore the principal-part quotient.  Its map to the full de Rham
complex is the boundary map
$P_V(J_0)\to R\Gamma(X,J_0)[1]$ of the meromorphic resolution.

This is an identification of crystals, including across collisions.
Under two lifts of a tuple agreeing on reduction, localize the same
bundle by the two divisor equations.  Their supports agree on reduction,
so their localizations agree after passing to all powers of those
equations.  These cofinal pole systems give the principal-part comparison;
in local coordinates it is differentiation in the marking variable with
the curve variable fixed.  Thus the support calculation gives the
specified comparison, rather than an identification only on geometric
fibers.

For $r$ factors, use $r$ curve variables $z_1,\ldots,z_r$ and the incidence
conditions $z_j\in V$.  ! base change and K\"unneth identify the tensor
product in \eqref{eq:T18} with the corresponding product support
calculation.  No smoothness of the intersections of these incidences is
needed: each support operation is defined by its localization triangle.
Project this incidence space to $X^r$.  Its relative Ran direction consists
of finite subsets containing both $U$ and all the specified $z_j$.
Universal homological contractibility identifies integration in this
direction with the identity, by the Ran trace
\cite[Theorem~1.6.5, Section~2.5.12, and Corollary~2.5.13]
{GaitsgoryContractibility}.  The resulting morphism is precisely the
one obtained by forgetting the supports in \eqref{eq:T17}, so it is the
map displayed in \eqref{eq:T18}.

All operations commute with the finite permutation actions and with
shifts.  In characteristic zero, symmetric powers are direct summands of
tensor powers.  The same calculation therefore applies to every symmetric
power, and continuity permits their direct sum.
\end{proof}

\subsection{The unipotent calculation}

We begin by identifying the period map on induced generators.  At a tuple
$V$, put $H=H_{N,V}^{\mathrm{lv}}$ and $h=\operatorname{Lie}H$.  Let $W$ be
an inducing representation, and let
$\operatorname{ev}_V:B_N\to BH$ classify the frame torsor
$\widetilde B_N\to B_N$.  The symbols $B_N$ and $\widetilde B_N$ here
include the fixed level data, as in Section~\ref{sec:unipotent}.

\begin{lemma}\label{lem:ran-shapiro}
With continuous cochains and the Harish--Chandra convention of
\eqref{eq:T13}, there are natural identifications
\begin{equation}\label{eq:T19}
\begin{aligned}
 B_V^{\mathrm{lv}}(\operatorname{Ind}W)
   &\simeq C^\bullet(h,W),\\
 (l^{\mathrm{lv}})^{-1}[-\delta]\otimes
       \operatorname{Per}_{N,V}^{\mathrm{lv}}(\operatorname{Ind}W)
   &\simeq C^\bullet\bigl(h,W\otimes
                                  \mathcal O(\widetilde B_N)\bigr)\\
   &\simeq R\Gamma\bigl(B_N,\operatorname{ev}_V^*W\bigr).
\end{aligned}
\end{equation}
The map \eqref{eq:T10} becomes the cochain pullback induced by
$W\to W\otimes\mathcal O(\widetilde B_N)$.
\end{lemma}

\begin{proof}
Use the increasing compact Lie lattices $A_l$ from \eqref{eq:T14}.
For a fixed $l$, a finite stage of $W$, and, when present, a finite stage
of functions, choose a sufficiently deep normal quotient $J$ through
which the requisite actions factor.  Finite-dimensional Lie Shapiro then
computes induction to $A_l/J$.  In cochain form its map inserts the
inducing vectors and wedges with the top exterior power of the
complementary modes.  The latter is exactly the determinant conversion
in \eqref{eq:T14}.  Thus the comparison is the map of Lie chains defining
\eqref{eq:T10}; with functions included, the projection formula for Lie
induction gives the same map.  The character $-\chi$ restricts to zero on
$h$, so it does not change these finite Shapiro computations.

There are two compatible exhaustions.  Deepening $J$ gives cochain
inflation, and increasing $A_l$ gives the same complementary top wedge as
in the semi-infinite transition maps.  Every inducing-module element
occurs for some $l$.  For any chosen finite representation and function
stages, the quotient can then be taken sufficiently deep and normal for
all of those data.  Taking their filtered colimit proves the first two
identifications of \eqref{eq:T19}, together with the asserted map.  In
the function calculation the stages are
$\overline Z=\widetilde B_N/J$ from \eqref{eq:T11}.

In characteristic zero, algebraic cohomology of a split unipotent group
is computed by its Lie cochains, also relatively over the present base.
For completeness, resolve a representation by cofree representations of
the form functions tensored with a representation.  Their Lie cochain
complex is the algebraic de Rham complex of the unipotent group, hence is
acyclic in positive degrees; successive vector-group extensions give
the relative assertion.  Apply this at the finite quotients and then
pass to the displayed colimit.  Descent for the frame torsor gives the
last identification in \eqref{eq:T19}.

These identifications also compute the crystals in the added markings.
Across a nilpotent lift comparison, the arcs and the frame torsor use the
same formal curves.  The finite lattice bounds have common enlargements,
and the quotient levels have common refinements, so the preceding
cochain maps compare on a cofinal system.  Evaluation is the map of those
same frame torsors.  The resulting comparison is therefore the
parameter comparison, not merely the underlying pointwise Shapiro map.
\end{proof}

Since $H$ is pro-unipotent, the induced trivial representation suffices
for the generation test of \eqref{eq:L2}.  Proposition
\ref{prop:ran-integrated-period} for $N$ consequently reduces to proving
that evaluation induces an isomorphism
\begin{equation}\label{eq:T20}
 \int_{V\supseteq U}
      C^\bullet\bigl(\operatorname{Lie}H_{N,V}^{\mathrm{lv}},
                    \mathcal O\bigr)
 \longrightarrow R\Gamma(B_N,\mathcal O).
\end{equation}
Indeed, the right side in \eqref{eq:T19} is the constant family in the
added markings.  Propagation \eqref{eq:T3} identifies it with the family
from the old tuple, and integration of that ! pull is the identity.

\begin{lemma}\label{lem:ran-unipotent-cohomology}
The evaluation map \eqref{eq:T20} is an isomorphism, with its relative
parameter comparisons.
\end{lemma}

\begin{proof}
First suppose that the group is the additive group of a vector bundle
$J$ on the curve.  Residue duality identifies continuous arc cochains
with
\[
 \operatorname{Sym}\bigl(P_V(J^*\otimes\Omega_X)[-1]\bigr),
\]
where tensor operations use the ! family convention.  A two-term perfect
complex for $R\Gamma(X,J)$ presents the stack of $J$-torsors as a vector
stack.  Computing its functions by additive-group cochains gives
\[
 R\Gamma(B_J,\mathcal O)
   \simeq\operatorname{Sym}\bigl(R\Gamma(X,J)^*[-1]\bigr).
\]
Serre duality identifies $R\Gamma(X,J)^*[-1]$ with
$R\Gamma(X,J^*\otimes\Omega_X)$.  Thus Lemma
\ref{lem:ran-principal-parts}, applied degree by degree and with the
shift $[-1]$, identifies these two symmetric algebras after integration.

It also identifies their actual comparison map.  Evaluation at a finite
jet is restriction of linear complexes,
\[
 R\Gamma(X,J)\longrightarrow\Gamma(X,J|_{V_{\mathrm{thick}}}).
\]
The shifted dual is its map on cochain generators.  Serre duality and
the lattice/support exact sequence identify this dual map with forgetting
supports in \eqref{eq:T17}.  Passing to all thickened divisors gives the
full principal-part map.  Over an infinitesimal test on which two tuples
agree on reduction, their thickenings are cofinal; the restriction maps
and the dual support maps consequently give the same crystal comparison.
This checks the derived map before truncating the linear complexes, and
proves \eqref{eq:T20} for the vector group.

We pass from the vector group to $N$ by a deformation whose gradings
control specialization.  In exponential coordinates replace the bracket
on $\mathfrak n^{\mathrm{lv}}$ by $t[-,-]$, over $\mathbb A^1_t$.
The positive root-height decomposition makes this construction global.
At $t=0$ the group is the additive group of
$J=\mathfrak n^{\mathrm{lv}}$, and at $t=1$ it is the original group.
The deformation concerns only the untwisted cochains and functions in
\eqref{eq:T20}.  Both sides and its evaluation map are computed
quasi-coherently in $t$, while retaining the relative D-module operations
in the added marking schemes.

These computations commute with specialization.  On the source, the
continuous Chevalley--Eilenberg complexes use locally free discrete duals
and their filtered unions.  On the target, use the presentations
\eqref{eq:T11}, taking the root-lattice enlargement sufficiently large
that $H^1$ vanishes in each of the finitely many root pieces.  The same
choice works over $\mathbb A^1_t$, since these underlying vector bundles
are unchanged.  The height lattices, their global-section subgroups
$K_l$, and the deep congruence quotients are invariant under both scalings
used below, by their exponential-coordinate definitions.  Finite quotient
presentations and their cochain computations are flat over the $t$-line.  Common deeper quotients compare
them with the exhausted presentation.  Evaluation uses these very
presentations, so its map commutes with specialization as well.

Here is the bound that makes the special fiber decisive.  Give a
coordinate, or a linear cochain, belonging to a height-$m$ root piece
first weight $m>0$.  Give $t$ first weight zero.  A second action is
\[
             t\longmapsto c t,\qquad x\longmapsto c^{-1}x
             \quad\text{in exponential coordinates}.
\]
It acts by group isomorphisms in the family: scaling the variables by
$c^{-1}$ compensates for scaling the bracket parameter by $c$.
Thus a linear coordinate or cochain has second weight $-1$, and $t$ has
second weight $1$.  In first weight $n\geq0$, every monomial has at most
$n$ such linear factors, and hence has second weight at least $-n$.
This bound applies to arc cochains and to the function modules in the
finite unipotent presentations.  In particular, functions on
$K_l\backslash(L_l/J)$ form a graded submodule of
$\mathcal O(L_l/J)$.  The bound therefore passes to these quotient
function modules, to the cochain complexes computing the quotient
stacks, and to their cohomology.  At finite stages relative unipotent cohomology may be
computed by finite-range Lie cochains, so these are actual graded
complexes.  The first grading is a direct sum, not a completion.

All maps in the added marking variables preserve both gradings.  Relative
de Rham pushforward, followed by the Ran colimits, preserves their direct
sums and the lower bound $-n$ in every first weight.  Let $Q_n$ be the
first-weight-$n$ part of the cone of \eqref{eq:T20} for the family.  The
vector-group calculation says
$Q_n\otimes^{\mathbf L}_{\mathcal O[t]}\mathcal O=0$ at $t=0$.
The associated exact triangle says that multiplication by $t$ is an
isomorphism on $Q_n$.  On cohomology it would therefore give isomorphisms
between successive second weights.  A nonzero class could be pulled back
repeatedly to weights below $-n$, which is impossible.  Hence every $Q_n$
vanishes.  Their direct sum is the full cone, so the family map, and in
particular its specialization at $t=1$, is an isomorphism.
\end{proof}

Together with Lemma~\ref{lem:ran-shapiro}, this proves Proposition
\ref{prop:ran-integrated-period} for $N$.  The factor
$l^{\mathrm{lv}}[\delta]$ is constant in the added markings and can be
tensored through the calculation.  We next compare restriction from $G$
with localization; this supplies the passage to the reductive group.

\subsection{Restriction and the two insertion comparisons}

For $G$, start at $U=T$ with a nonzero level covering the original tuple,
as in the congruence calculation of Section~\ref{sec:localization}.
Fix any such level.  The additional depth required by a finite
Grassmannian kernel will enter only when we apply \eqref{eq:T6} below.
Combining \eqref{eq:T5}, propagation, and trace gives a map
\begin{equation}\label{eq:T21}
 \int_{U\supseteq T}
  \operatorname{Loc}_{N,U}^{\mathrm{lv}}
       \bigl(\operatorname{Res}\operatorname{ins}_{G,T,U}M\bigr)
 \longrightarrow p^!\operatorname{Loc}_{G,T}^{\mathrm{lv}}(M).
\end{equation}

The reductive comparison fits into the following square.  To display its
four vertices without repeating the insertion expressions, set
\[
 I_U=\operatorname{Res}\operatorname{ins}_{G,T,U}M,
 \qquad c_N=l_N[-\delta_N],
\]
and define
\[
\begin{aligned}
 \mathsf B_2(M)
   &=c_N\otimes\int_{T\subseteq U\subseteq V}
       B_V^0\bigl(\operatorname{ins}_{N,U,V}I_U\bigr),\\
 \mathsf B_1(M)
   &=c_N\otimes\int_{T\subseteq V}B_V^0(I_V),\\
 \mathsf P_N(M)
   &=\int_{T\subseteq U}\operatorname{Per}_{N,U}^{\mathrm{lv}}(I_U)
                                            [-2\delta],\\
 \mathsf P_G(M)
   &=\operatorname{Per}_{G,T}^{\mathrm{lv}}(M)[-2\delta].
\end{aligned}
\]
The level remains on $T$, so $\delta$ and $c_N$ are constant in the added
markings.  The maps form the commutative diagram
\begin{equation}\label{eq:ran-reductive-square}
\begin{tikzcd}[column sep=8em,row sep=4em]
 \mathsf B_2(M)
   \arrow[r,"\text{integrated $N$ case of \eqref{eq:T16}}","\sim"']
   \arrow[d,"\text{insertion comparison}"']
 &\mathsf P_N(M)
   \arrow[d,"\text{period of \eqref{eq:T21}}"]\\
 \mathsf B_1(M)
   \arrow[r,"\text{$G$ case of \eqref{eq:T16}}"']
 &\mathsf P_G(M).
\end{tikzcd}
\end{equation}
The top map is already an isomorphism by the unipotent calculation.  The
right map applies the period operation to \eqref{eq:T21}.  For the left
map, use the natural insertion comparison
\[
 \operatorname{ins}_{N,U,V}I_U\longrightarrow I_V
\]
and the integration counit for forgetting $U$.  By \eqref{eq:T9}, this is
the map \eqref{eq:T23} in the full-arc normalization displayed here.
The square commutes because \eqref{eq:T10} is natural, the restriction
map \eqref{eq:T5} commutes with propagation \eqref{eq:T3}, and the
integration traces compose; the $G$-period map uses exactly the resulting
restriction map \eqref{eq:T15}.  We will prove that the right vertical
map is an isomorphism using \eqref{eq:T21}, and then that the left
vertical map is an isomorphism using \eqref{eq:T23}.  This will prove the
bottom map is an isomorphism.

\begin{lemma}\label{lem:ran-restriction-localization}
The map \eqref{eq:T21} is an isomorphism.
\end{lemma}

\begin{proof}
Pull to bundle charts on $B_N$ and apply $b$ in the old tuple
variables.  ! base change commutes with the indicated integrations.  For this
isomorphism test we may forget differential operators in the bundle-chart
variables, while retaining the D-module operations in the added markings.
By pro-unipotence at the initial congruence level and \eqref{eq:L2}, it
suffices to use the arc-induced generator with trivial inducing
representation.  At a chart bundle $P$ with unipotent data, filter the
calculation \eqref{eq:T2} by total PBW degree, including the Lie-chain
factors.  The associated graded on the left, before integration, is
\begin{equation}\label{eq:T22}
 \operatorname{Sym}\!\left(
   \operatorname{Cone}\!\left(
       \mathfrak n^{\mathrm{lv}}_{P,\mathrm{out},U}
          \longrightarrow P_U(\mathfrak g^{\mathrm{lv}}_P)
                         \right)\right).
\end{equation}
On the right the out-current term is
$\mathfrak g^{\mathrm{lv}}_{P,\mathrm{out},U}$ instead.  Its cone is
$R\Gamma(X,\mathfrak g^{\mathrm{lv}}_P)[1]$ by the meromorphic resolution.
The comparison is induced by inclusion of out currents.

To apply integration to this filtration, we check it in the added
markings.  The bundle $P$ and its level are fixed over these variables.
Across nilpotent lifts of the added markings, its disk frames are frames
on the same formal curves.  Changes of frame belong to the regular level
arc group and preserve the PBW filtration by Harish--Chandra equivariance.
The out-current and principal-part complexes use these same comparisons.
Hence \eqref{eq:T22} is a filtered crystal calculation in the added
markings.  There is no assertion here that this is a filtration by
D-submodules in the bundle-chart variables; those variables have already
been forgotten for the isomorphism test.

Set $J_P=\mathfrak g^{\mathrm{lv}}_P/\mathfrak n^{\mathrm{lv}}_P$.
The smooth root lattices make
$0\to\mathfrak n^{\mathrm{lv}}_P\to
\mathfrak g^{\mathrm{lv}}_P\to J_P\to0$ an exact sequence of vector
bundles.  Affineness of the complement of $U$ and the principal-part
sequence give a diagram of exact triangles.  Its common first term is
$K_P=R\Gamma(X,\mathfrak n^{\mathrm{lv}}_P)[1]$; its middle terms are the
two cone-generating complexes just described; and its quotient map is
\[
        P_U(J_P)\longrightarrow R\Gamma(X,J_P)[1].
\]
More explicitly, the two rows of the diagram are
\[
\begin{array}{cccccc}
 K_P&\longrightarrow&
 \operatorname{Cone}(\mathfrak n^{\mathrm{lv}}_{P,\mathrm{out},U}
                  \to P_U(\mathfrak g^{\mathrm{lv}}_P))
 &\longrightarrow&P_U(J_P)&\longrightarrow K_P[1],\\[2pt]
 K_P&\longrightarrow&R\Gamma(X,\mathfrak g^{\mathrm{lv}}_P)[1]
 &\longrightarrow&R\Gamma(X,J_P)[1]&\longrightarrow K_P[1],
\end{array}
\]
with identity on $K_P$ and the forget-support map on the quotients.
For each symmetric degree $r$, the extension filtration has $r+1$ steps,
with pieces given by tensor products of a symmetric power of $K_P$ and a
symmetric power of the corresponding quotient.  The former is constant
in $U$; Lemma~\ref{lem:ran-principal-parts} identifies the integrated maps
on the latter.  Thus the comparison is an isomorphism in each symmetric
degree.  Finite induction through the extension filtration, and then the
exhaustive PBW colimit, prove \eqref{eq:T21}.
\end{proof}

Lemma~\ref{lem:ran-restriction-localization}, followed by the period
operation, proves that the right vertical map in
\eqref{eq:ran-reductive-square} is an isomorphism.  It remains to prove
the same assertion for the left vertical map.  By \eqref{eq:T9}, we may
remove the common invertible normalization and use the level-lattice
form of this comparison.  Thus we must show that
\begin{equation}\label{eq:T23}
\begin{aligned}
 &\int_{T\subseteq U\subseteq V}
 B_V^{\mathrm{lv}}\bigl(
    \operatorname{ins}_{N,U,V}\operatorname{Res}
                    \operatorname{ins}_{G,T,U}M\bigr)\\
 &\hspace{30mm}\longrightarrow
 \int_{T\subseteq V}B_V^{\mathrm{lv}}\bigl(
                 \operatorname{Res}\operatorname{ins}_{G,T,V}M\bigr)
\end{aligned}
\end{equation}
is an isomorphism.  Both nested integrals use the ordinary relative
parameter operations, with ! pulls and their integration counits.

\begin{lemma}\label{lem:ran-two-insertions}
The map \eqref{eq:T23} is an isomorphism.
\end{lemma}

\begin{proof}
Again use the trivially induced generator at the old congruence level.
Filter restriction by the number of outside modes as in \eqref{eq:T4}.
Insertion \eqref{eq:L9}, followed by Lemma~\ref{lem:ran-shapiro}, identifies
the map in outside degree $i$ with the integrated map of complexes
\begin{equation}\label{eq:T24}
\begin{aligned}
 C^\bullet\!\left(h_V,
       \operatorname{Sym}^{i}P_U
          (\mathfrak g^{\mathrm{lv}}/\mathfrak n^{\mathrm{lv}})\right)
 \longrightarrow
 C^\bullet\!\left(h_V,
       \operatorname{Sym}^{i}P_V
          (\mathfrak g^{\mathrm{lv}}/\mathfrak n^{\mathrm{lv}})\right).
\end{aligned}
\end{equation}
Here $h_V=\operatorname{Lie}H_{N,V}^{\mathrm{lv}}$.  The map on the
coefficients is the actual principal-part inclusion: in the inserting
induction, the intermediate Laurent vectors have poles only on $U$.
For a fixed outside degree and finite pole bound the inducing
representations in \eqref{eq:T4} are finite stages in heart degree; choose
a common sufficiently deep quotient, apply Shapiro there, and pass to the
filtered union of these stages.  This gives \eqref{eq:T24}, including its
map.  The old level remains fixed, so all these filtrations use the
formal-curve comparisons in the added variables already checked for
\eqref{eq:T22}.

We reduce \eqref{eq:T24} to tensor products of principal parts by two
finite filtrations.  First fix $i$.  The heights of the root and Cartan
pieces of $\mathfrak g/\mathfrak n$ lie in a finite interval, say
$[a,b]$; a degree-$i$ symmetric monomial has height in $[ia,ib]$.
The arc-unipotent action strictly raises height.  The resulting finite
filtration, of length at most $i(b-a)+1$, has trivial action on its
associated graded.  Thus each piece is the tensor product of
$C^\bullet(h_V,\mathcal O)$ with a fixed tensor/symmetry expression in
coefficient principal parts.

Next decompose the plain cochain complex by the positive height grading
used in Lemma~\ref{lem:ran-unipotent-cohomology}.  In first weight $n$,
cochain degree is at most $n$, since every linear cochain has positive
height at least one.  Its cochain-degree filtration is therefore finite.
By residue duality each of its terms is a tensor, shift, and alternating
summand of the principal parts of the dual root bundles with
$\Omega_X$.  Consequently the comparison on every such term is a
comparison between a finite list of principal parts taken at $U$ and at
$V$, on the left, and the same list all taken at $V$, on the right.

For these terms use the incidence proof of \eqref{eq:T18}.  Both sides
project to the same power of the curve.  On the left, $U$ must contain
$T$ and the curve variables assigned to its list; $V$ must contain $U$
and the remaining variables.  Integrating first in the larger tuple and
then in the smaller tuple contracts these two Ran directions.  On the
right, a single tuple $V$ contains $T$ and all the curve variables, and
contracts by the same trace.  The comparison is the isomorphism between
these two contractions: the inclusion of incidences $z\in U\Rightarrow
z\in V$ uses exactly the forget-support maps of \eqref{eq:T17}, whose
traces compose.

Finite induction through the cochain-degree and coefficient-height
filtrations proves the assertion in each first weight and outside
symmetric degree.  Take the direct sum of first weights and then the
exhaustive outside-mode filtration.  Continuity of the cochain,
insertion, and integration constructions permits these colimits.
This proves \eqref{eq:T23}.
\end{proof}

The two insertion comparisons finish the reductive case of Proposition
\ref{prop:ran-integrated-period}.  Every comparison was made in the
parameter operations, with common deeper levels on any finite collection
of bounds.  It follows by continuous extension that the resulting
isomorphism has all the compatibilities stated there.

\subsection{Normalization and descent from Grassmannian tests}

We now complete the proof of Theorem~\ref{loc:period-comparison}.
Start with a test
$W'=\mathbf d(K)$, where $K$ lies in a finite closed Grassmannian bound.
Choose the sufficiently deep level required by \eqref{eq:T6}.  Formula
\eqref{eq:T8} expresses the unnormalized pairing as a shifted period.
Proposition~\ref{prop:ran-integrated-period}, together with the
normalization of localization, cancels $l_N[-\delta_N]$ and gives
\begin{equation}\label{eq:T25}
       \int_{V\supseteq T}
       B^0_V\bigl(\operatorname{ins}_{T,V}(K\star M)\bigr).
\end{equation}
In particular all auxiliary level determinants and dimension shifts have
canceled; the semi-infinite lattice remaining in this formula is the full
arc lattice $h_V^0$.

The insertion-with-action assertion of \eqref{eq:T6} and the descriptions
\eqref{eq:L3} and \eqref{eq:L6} identify \eqref{eq:T25} with the right
side of \eqref{eq:T1}.  More precisely, compare $B^0$ on the
Harish--Chandra categories with the Whittaker coinvariant projection
after their forgetful maps to $\mathcal M$, followed by
$\operatorname{oblv}_W$.  After applying
$b$, the forgetful maps to $\mathcal M^0$ preserve compact objects by
the finite relative Chevalley calculation of
Proposition~\ref{loop:hc-renormalization}.  Both operations on these
compact presentations are the BRST complex \eqref{eq:L6} with lattice
$h_V^0$; continuity then gives their agreement on all inputs.  Insertion
of $\mathbf d(K)$ is the closed Grassmannian image from the tests version
of \eqref{eq:L9}.  Thus the comparison is a comparison in parameters,
including moving kernels.  Every operation used here is the finite-bound
pairing followed by continuous extension.

We must still descend the comparison from these tests to all of $d_T$.
The construction is compatible as the bounds grow: for any finite set of
bounds, take one common sufficiently deep level, and use the pull and
period comparisons when deepening further.  We have therefore constructed
the comparison after precomposition by
\[
 \mathbf d:D_{\mathcal L^{-s}}(\operatorname{Gr}_T)\longrightarrow d_T.
\]
Precomposition along this functor is fully faithful on continuous
functionals.  Indeed, its transpose is the full embedding of the
positive-character Whittaker category on the Grassmannian, by
\eqref{loop:whittaker-coinvariants}--\eqref{eq:loop-whittaker-functional}
and the normalization
\eqref{eq:W2}.  After currying, maps between the two pairings are therefore
computed fully faithfully after precomposition by $\mathbf d$.
The isomorphism just constructed, and its inverse, descend uniquely.

For a pairing valued in D-modules on a base, apply the same argument after
pairing with the dual of the parameter D-module category.  It reduces the
claim to continuous Vect-valued pairings, to which the preceding full
faithfulness applies.  The argument also applies with any list of
parameter actions, so the descended isomorphism is parameter-linear and
respects the action maps.  This proves \eqref{eq:T1} with all of its
stated compatibilities.

\subsection{The Ran pairing and reversal of levels}

To pass to independent marking tuples, write $C=C_{\mathrm{Ran}}$ and
$\check C=\check C_{\mathrm{Ran}}$.  In this subsection only, $d=d_{\mathrm{Ran}}$
denotes the tests category, rather than the numerical shifted level.
The equivalences \eqref{eq:L11}, in particular
$C\simeq\check d$ and $\check C\simeq d$, extend to Ran by their diagonal
comparisons; the categories presented by ! systems also use the colimit
by closed-image operations.  Integrating normalized localization from
tuple schemes gives continuous functors
\[
                \Lambda:C\longrightarrow A,
       \qquad \check\Lambda:\check C\longrightarrow B'.
\]
Integration here has the stack and sheaf conventions of the earlier
parameter operations; on $X^I$ it is ordinary proper de Rham pushforward.

\begin{proposition}\label{prop:ran-pairing}
For $M\in C$ and $\check M\in\check C$, there is a canonical isomorphism
of bilinear forms
\begin{equation}\label{eq:T26}
 \bigl\langle R_A\Lambda(M),\check\alpha(\check M)\bigr\rangle
 \simeq
 \bigl\langle R_{B'}\check\Lambda(\check M),\alpha(M)\bigr\rangle.
\end{equation}
It is compatible with parameter actions, permutations, and diagonals.
\end{proposition}

\begin{proof}
Represent the inputs over independent marking schemes.  On their product,
insert each list into the union and apply \eqref{eq:T1}; then integrate
over the two original marking schemes.  We check that the resulting
pairings are precisely the displayed ones.

For Whittaker coefficients, use \eqref{eq:W2} with the tests extended by
closed push from the old pole data, where no extra poles have yet been
allowed.  ! restriction of the coefficient to these old pole bounds is
its old value.  This is the closed parameter/pole comparison established
in the tests-insertion calculation after \eqref{eq:L9}; equivalently,
compute \eqref{eq:W1} away from the union of marking points, where the
auxiliary correspondences are unchanged.  On the localization side use
propagation \eqref{eq:T3} with zero level divisor at the additional
markings.  ! projection formula and base change, followed by de Rham
integration, consequently compute the first pairing in \eqref{eq:T26}.
On the other side of \eqref{eq:T1}, the identification \eqref{eq:L12}
and the common underlying-complex functor \eqref{eq:walg-forgetful}
give the second one.

These are maps on the actual marking schemes.  Their permutation and
diagonal compatibilities were retained in the parameter constructions;
the level choices can be deepened to common choices in each comparison
diagram.  The same argument works for closed pushes in tuple coordinates,
using the parameter-linear calculation and the same closed-embedding
adjunction maps.  These diagrams present the Ran categories and their
continuous functors, so the isomorphism descends to Ran as asserted.
\end{proof}

Finally, the period calculation applies with simultaneous reversal of the
characters: take negative character in \eqref{eq:L1}--\eqref{eq:L3} and
\eqref{eq:T1} on their original-twist categories, and positive character
on their opposite categories.  It also applies to any other pair of
irrational inverse shifted forms.  The proof used no order or positivity
condition on those forms.

For negated shifted forms, one must first recompute the affine levels:
the affine level is the new shifted form minus the dual Coxeter number.
The determinant exponents are then exactly those prescribed in
\eqref{eq:L1}.  One may subsequently transport the comparison by integral
powers of the determinant.  These tensor identifications respect the
formulas: the pole, Grassmannian, and Hecke maps carry their relative
determinant; the unipotent bundle and action maps use the filtration
trivialization and orbit maps of \eqref{eq:W1}; and the two powers cancel
in opposite pairings.  This is the sign convention used in
Section~\ref{sec:equivalence}.
